如上所述, 建立存在性理论的关键是证明类似于 Gårding 不等式的能量估计.

\begin{Theorem}
\label{thm:neumann-energy-estimate}
设 $L(\bm{x},D)$ 为 \eqref{eq:second-order-garding-divergence-operator} 所述形式的二阶散度形式线性偏微分算子, 且以某个 $\theta>0$ 满足一致椭圆性条件 \eqref{eq:second-order-ellipticity-negative-coefficient-matrix}. 再假设 $a_{ij},b_k\in L^\infty(\Omega)$, 其中 $i,j=1,\ldots,n$, $k=0,\ldots,n$. 那么存在常数 $\overline c$ 和 $\lambda_N\geq0$, 使得
\begin{equation}\label{eq:neumann-energy-coercivity-inequality}
B[u,u]+\lambda_N\|u\|_{L^2(\Omega)}^2\geq\overline c\|u\|_{H^1(\Omega)}^2\quad\text{for all }u\in H^1(\Omega),
\end{equation}
其中 $B$ 由 \eqref{eq:second-order-garding-bilinear-form} 定义.
\end{Theorem}

\begin{Proof}
该定理的陈述与 Theorem~\ref{thm:second-order-garding-inequality} 相同, 只是现在要在空间 $H^1(\Omega)$ 上证明, 而不再仅限于 $H_0^1(\Omega)$. 因而证明中唯一的差别是不能再使用 Poincaré 不等式. 然而, 仍可如前得到
\begin{equation}\label{eq:neumann-energy-gradient-lower-bound}
B[u,u]\geq\dfrac{\theta}{2}\int_\Omega|\nabla u|^2\,\mathrm{d}\bm{x}-C\|u\|_2^2.
\end{equation}
此时不用 Poincaré 不等式, 而只须写为
\begin{equation}\label{eq:neumann-energy-full-h1-lower-bound}
B[u,u]\geq\dfrac{\theta}{2}\|u\|_{1,2}^2-\left(C+\dfrac{\theta}{2}\right)\|u\|_2^2,
\end{equation}
证明即告完成.
\end{Proof}

值得注意的是, 在 Theorem~\ref{thm:second-order-garding-inequality} 的证明中使用 Poincaré 的唯一真正原因, 是得到关于常数 $\lambda_G$ 的更精确估计. 然而, 既然本来就没有试图明确给出最优常数, 这一番努力多少有些白费.

有了满足 Lax-Milgram 引理强制性条件的能量估计, 只须对 Theorem~\ref{thm:weak-dirichlet-positive-shift-existence} 的证明稍加修改, 就得到 Neumann 问题弱解的存在性.

\begin{Theorem}
\label{thm:neumann-positive-shift-weak-existence}
设 $L(\bm{x},D)$ 为二阶散度形式线性偏微分算子, 满足 Theorem~\ref{thm:neumann-energy-estimate} 的假设. 那么存在 $\lambda_N\geq0$, 使得对任意 $\widetilde\lambda\geq\lambda_N$ 和任意 $f\in L^2(\Omega)$, 算子
\begin{equation}\label{eq:neumann-positive-shifted-operator}
\widetilde L(\bm{x},D):=L(\bm{x},D)+\widetilde\lambda
\end{equation}
的 Neumann 问题存在唯一弱解 $u\in H^1(\Omega)$. 此外, 该解满足
\begin{equation}\label{eq:neumann-weak-solution-l2-data-bound}
\|u\|_{1,2}\leq C\|f\|_2.
\end{equation}

\end{Theorem}
